\setcounter{section}{23}

  \zwischenueberschrift{Übungsaufgaben}

\inputaufgabe
{}
{

Diskutiere die
[[Riemann-Integral/Hauptsatz/Newton-Leibniz/Fakt|Newton-Leibniz-Formel]]
als einen Spezialfall des
[[Mannigfaltigkeit mit Rand/Satz von Stokes/Fakt|Satzes von Stokes]].

}
{} {}

\inputaufgabe
{}
{

Beweise die
[[Satz von Stokes/Quaderversion/Fakt|Quaderversion des Satzes von Stokes]]
direkt für einen Quader
\mathl{[a_1,b_1] \times \cdots \times [a_n,b_n]}{} und eine
$n-1$-\definitionsverweis {Differentialform}{}{}
$\omega$ der Gestalt

\mavergleichskettedisp
{\vergleichskette
{ \omega
}
{ =} { x_1^{m_1 } \cdots x_n^{m_n }  dx_2 \wedge \ldots \wedge dx_n
}
{ } {
}
{ } {
}
{ } {
}
}
{}{}{.}

}
{} {}

\inputaufgabe
{}
{

Es sei $M$ eine
$n$-\definitionsverweis {dimensionale}{}{}
\definitionsverweis {orientierte}{}{}
\definitionsverweis {differenzierbare Mannigfaltigkeit mit Rand}{}{}
und mit
\definitionsverweis {abzählbarer Basis der Topologie}{}{}
und es sei $\omega$ eine
\definitionsverweis {stetig differenzierbare}{}{} \definitionsverweis {geschlossene}{}{} $(n-1)$-\definitionsverweis {Differentialform
}{}{}
auf $M$ mit
\definitionsverweis {kompaktem Träger}{}{.}
Zeige

\mavergleichskettedisp
{\vergleichskette
{ \int_{\partial M} \omega
}
{ =} { 0
}
{ } {
}
{ } {
}
{ } {
}
}
{}{}{.}

}
{} {}

\inputaufgabe
{}
{

Es sei $M$ eine
\definitionsverweis {kompakte}{}{}
$n$-\definitionsverweis {dimensionale}{}{}
\definitionsverweis {orientierte}{}{}
\definitionsverweis {differenzierbare Mannigfaltigkeit}{}{}
\zusatzklammer {ohne Rand} {} {}
mit
\definitionsverweis {abzählbarer Basis der Topologie}{}{}
und es sei $\omega$ eine
\definitionsverweis {stetig differenzierbare}{}{}
$(n-1)$-\definitionsverweis {Differentialform
}{}{}
auf $M$. Zeige

\mavergleichskettedisp
{\vergleichskette
{
\int_{ M  }  d\omega
}
{ =} { 0
}
{ } {
}
{ } {
}
{ } {
}
}
{}{}{.}

}
{Was bedeutet diese Aussage für $S^1$? Wie kann man diese Aussage in diesem Fall über ein Wegintegral beweisen?} {}

\inputaufgabe
{}
{

Es sei $M$ eine
\definitionsverweis {kompakte}{}{}
\definitionsverweis {orientierte}{}{}
\definitionsverweis {differenzierbare Mannigfaltigkeit}{}{}
\zusatzklammer {ohne Rand} {} {}
mit
\definitionsverweis {abzählbarer Basis der Topologie}{}{}
und es sei $\tau$ eine
\definitionsverweis {positive Volumenform}{}{}
auf $M$. Zeige, dass $\tau$ nicht
\definitionsverweis {exakt}{}{}
ist.

}
{Wie sieht dies ohne die Kompaktheitsvoraussetzung aus?} {}

\inputaufgabegibtloesung
{}
{

Es sei $M$ eine
\definitionsverweis {differenzierbare Mannigfaltigkeit}{}{}
und $\omega$ eine
\definitionsverweis {exakte}{}{} \definitionsverweis {Differentialform}{}{}
auf $M$. Es sei $S$ eine
\definitionsverweis {kompakte}{}{}
\definitionsverweis {orientierte}{}{}
\definitionsverweis {differenzierbare Mannigfaltigkeit}{}{}
\zusatzklammer {ohne Rand} {} {}
mit
\definitionsverweis {abzählbarer Basis der Topologie}{}{}
und es sei
\maabbdisp {\varphi} { S } { M
} {}
eine
\definitionsverweis {stetig differenzierbare}{}{}
Abbildung. Zeige

\mavergleichskettedisp
{\vergleichskette
{
\int_{  S  }   \varphi^*\omega
}
{ =} { 0
}
{ } {
}
{ } {
}
{ } {
}
}
{}{}{.}

}
{} {}

\inputaufgabe
{}
{

Wir betrachten die
\definitionsverweis {stetig differenzierbare}{}{}
Abbildung
\zusatzklammer {\mathlk{n \geq 2}{}} {} {}
\maabbeledisp {\pi} {\R^n \setminus \{0\}} { S^{n-1}
} {(x_1 , \ldots , x_n) } { { \frac{   1  }{  \sqrt{x_1^2 + \cdots + x_n^2}  } } (x_1 , \ldots , x_n)
} {.}
Es sei $\omega$ die
\definitionsverweis {kanonische Volumenform}{}{}
auf $S^{n-1}$. Zeige, dass
\mathl{\pi^* \omega}{} auf $\R^n \setminus \{0\}$ eine
\definitionsverweis {geschlossene}{}{,}
aber keine
\definitionsverweis {exakte}{}{}
$n-1$-\definitionsverweis {Differentialform}{}{}
ist.

}
{} {}

\inputaufgabe
{}
{

Es sei

\mavergleichskette
{\vergleichskette
{ M
}
{ = }{ B  \left(  0 , 1 \right) \setminus \{(0,0)\}
}
{ \subset }{ \R^2
}
{    }{
}
{   }{
}
}
{}{}{}
mit dem Rand

\mavergleichskette
{\vergleichskette
{ \partial M
}
{ = }{ S^1
}
{  }{
}
{    }{
}
{   }{
}
}
{}{}{.}
Wir betrachten die beiden
\definitionsverweis {stetig differenzierbaren}{}{}
$1$-\definitionsverweis {Differentialformen}{}{}

\mathdisp {\omega= ydx-xdy \text{   und } \tau= { \frac{   1  }{  x^2+y^2 } } \left(  ydx-xdy  \right)} {  }

auf $M$. Zeige, dass die Einschränkungen der beiden Formen auf den Rand übereinstimmen und insbesondere

\mavergleichskette
{\vergleichskette
{ \int_{\partial M} \omega
}
{ = }{ \int_{\partial M} \tau
}
{  }{
}
{    }{
}
{   }{
}
}
{}{}{}
gilt. Vergleiche
\mathkor {} {\int_M d \omega} {und} {\int_M d \tau} {.}

}
{} {}

\inputaufgabe
{}
{

Man mache sich klar, dass der
[[Integration auf ebener Mannigfaltigkeit mit Rand/Satz von Green/Flächenversion/Fakt|Satz von Green]]
nicht behauptet, dass der Flächeninhalt eines umrandeten Gebiets im $\R^2$ nur von der Länge des Randes abhängt.

}
{} {}

\inputaufgabe
{}
{

Es sei $D$ das durch
$(0,2),\, (1,-1)$ und $(-2,-1)$
gegebene Dreieck und

\mavergleichskette
{\vergleichskette
{ \tau
}
{ = }{ x^2ydx \wedge dy
}
{  }{
}
{    }{
}
{   }{
}
}
{}{}{}
eine
$2$-\definitionsverweis {Differentialform}{}{} auf $D$. Finde eine
\definitionsverweis {Stammform}{}{}
für $\tau$ und berechne damit
\mathl{\int_{ D  }  \tau}{} durch ein Integral über dem Dreiecksrand.

}
{} {}

\inputaufgabe
{}
{

Bestätige
den [[Integration auf ebener Mannigfaltigkeit mit Rand/Satz von Green/Fakt|Satz von Green]]
für das Einheitsquadrat

\mavergleichskette
{\vergleichskette
{ T
}
{ = }{ [0,1] \times [0,1]
}
{  }{
}
{    }{
}
{   }{
}
}
{}{}{}
und die
\definitionsverweis {Differentialformen}{}{}

\mavergleichskettedisp
{\vergleichskette
{ \omega
}
{ =} { x^ay^b dx + x^cy^d dy
}
{ } {
}
{ } {
}
{ } {
}
}
{}{}{}
mit

\mavergleichskette
{\vergleichskette
{ a,b,c,d
}
{ \in }{ \N
}
{  }{
}
{    }{
}
{   }{
}
}
{}{}{}
durch explizite Berechnungen.

}
{} {}

\inputaufgabe
{}
{

Bestätige
den [[Integration auf ebener Mannigfaltigkeit mit Rand/Satz von Green/Fakt|Satz von Green]]
durch explizite Berechnungen für die Menge
\mathl{T=[-2,2] \times [-2,2] \setminus U \left(   0 , 1  \right)}{}
\zusatzklammer {also das zentrierte Quadrat der Seitenlänge $4$ ohne den offenen Einheitskreis} {} {}
und die
\definitionsverweis {Differentialform}{}{}

\mavergleichskette
{\vergleichskette
{ \omega
}
{ = }{(x-y)dx + xy dy
}
{  }{
}
{    }{
}
{   }{
}
}
{}{}{.}

}
{} {}

\inputaufgabegibtloesung
{}
{

Beweise den Satz von Green für ein Dreieck $D$ mit den Eckpunkten

\mavergleichskette
{\vergleichskette
{ P_1,P_2,P_3
}
{ \in }{ \R^2
}
{  }{
}
{    }{
}
{   }{
}
}
{}{}{}
und für die Differentialform
\mathl{xdy}{.}

}
{} {}

\inputaufgabe
{}
{

Es sei
\maabbdisp {f} { [a,b] } { \R_{+}
} {}
eine
\definitionsverweis {stetig differenzierbare Funktion}{}{.}
Wir fassen den
\definitionsverweis {Subgraphen}{}{}
als eine
\definitionsverweis {Mannigfaltigkeit mit Rand}{}{}
auf, wobei der Rand aus dem Graphen, dem Grundintervall und den beiden Seitenkanten, aber ohne die vier Eckpunkte besteht. Berechne den Flächeninhalt des Subgraphen mit den beiden
\definitionsverweis {Differentialformen}{}{}

\mathl{xdy}{} und
\mathl{ydx}{} über den Rand.

}
{} {}

\inputaufgabe
{}
{

Bestimme das Volumen der dreidimsionalen
\definitionsverweis {abgeschlossenen Einheitskugel}{}{}
durch ein geeignetes Flächenintegral über die
\definitionsverweis {Einheitssphäre}{}{}
$S^2$.

}
{} {}

\inputaufgabegibtloesung
{}
{

Wir betrachten die
$2$-\definitionsverweis {Differentialform}{}{}

\mavergleichskettedisp
{\vergleichskette
{ \omega
}
{ =} { x dy \wedge dz -y dx \wedge dz + z dx \wedge dy
}
{ } {
}
{ } {
}
{ } {
}
}
{}{}{}
auf der Einheitskugel
\mathl{B  \left(  0 , 1  \right)}{.}
\aufzaehlungdreiabc{Zeige, dass $d \omega$ das Dreifache der Standardvolumenform auf der Kugel ist.
}{Zeige, dass
\mathl{\omega\vert_{S^2}}{} die Standardflächenform auf der Einheitssphäre ist.
}{Berechne die Kugeloberfläche aus dem
[[Allgemeines Kugelvolumen/Mit Cavalieri-Prinzip/Beispiel|Kugelvolumen]]
mit dem
[[Mannigfaltigkeit mit Rand/Satz von Stokes/Fakt|Satz von Stokes]].
}

}
{} {}

\inputaufgabegibtloesung
{}
{

Es sei $\omega$ eine stetig differenzierbare
$n-1$-\definitionsverweis {Differentialform}{}{}
auf einer
\definitionsverweis {offenen Menge}{}{}

\mavergleichskette
{\vergleichskette
{ U
}
{ \subseteq }{ \R^n
}
{  }{
}
{    }{
}
{   }{
}
}
{}{}{}
und sei die
\definitionsverweis {äußere Ableitung}{}{}
gleich

\mavergleichskettedisp
{\vergleichskette
{d \omega
}
{ =} { fdx_1 \wedge \ldots \wedge dx_n
}
{ } {
}
{ } {
}
{ } {
}
}
{}{}{}
mit einer Funktion
\maabb {f} { U } {\R
} {.}
Zeige für

\mavergleichskette
{\vergleichskette
{ P
}
{ \in }{ U
}
{  }{
}
{    }{
}
{   }{
}
}
{}{}{}
die Gleichheit

\mavergleichskettedisp
{\vergleichskette
{ f(P)
}
{ =} { { \frac{   1  }{  \beta_n  } } \cdot \operatorname{lim}_{   \epsilon \rightarrow  0  } \, { \frac{   1  }{  \epsilon^n } } \int_{S^{n-1}(P, \epsilon)} \omega
}
{ } {
}
{ } {
}
{ } {
}
}
{}{}{,}
wobei $\beta_n$ das
[[Allgemeines Kugelvolumen/Mit Cavalieri-Prinzip/Beispiel|Volumen]]
der $n$-dimensionalen Einheitskugel bezeichnet.

}
{} {}

Für einen Spezialfall der vorstehenden Aussage siehe
[[Vektorfeld/R^2/Gegenableitungdifferenz/Limesformel/Aufgabe|Aufgabe 58.23 (Analysis (Osnabrück 2021-2023))]].

\inputaufgabe
{}
{

Es sei $M$ eine
\definitionsverweis {differenzierbare Mannigfaltigkeit}{}{}
mit einem nichtleeren Rand. Zeige, dass es eine
\definitionsverweis {differenzierbare Abbildung}{}{}
\maabb {} { M } { \partial M
} {}
gibt.

}
{} {}

\inputaufgabe
{}
{

Es sei

\mavergleichskette
{\vergleichskette
{ H
}
{ \subseteq }{ \R^n
}
{  }{
}
{    }{
}
{   }{
}
}
{}{}{}
\zusatzklammer {
\mavergleichskette
{\vergleichskette
{ n
}
{ \geq }{ 1
}
{  }{
}
{    }{
}
{   }{
}
}
{}{}{}} {} {}
ein
\definitionsverweis {Halbraum}{}{.}
Zeige, dass es eine
\definitionsverweis {differenzierbare Abbildung}{}{}
\maabbdisp {\varphi} { H } { \partial H
} {}
gibt, deren
\definitionsverweis {Einschränkung}{}{}
auf
\mathl{\partial H}{} die
\definitionsverweis {Identität}{}{}
ist.

}
{Wie sieht das bei

\mavergleichskette
{\vergleichskette
{ n
}
{ = }{ 0
}
{  }{
}
{    }{
}
{   }{
}
}
{}{}{}
aus?} {}

\inputaufgabe
{}
{

Wir betrachten die
\definitionsverweis {Mannigfaltigkeit mit Rand}{}{}

\mavergleichskette
{\vergleichskette
{ M
}
{ = }{ \R^n \setminus  U \left(   0 , 1  \right)
}
{  }{
}
{    }{
}
{   }{
}
}
{}{}{.}
Zeige, dass es eine stetig differenzierbare Abbildung von $M$ auf den Rand
\mathl{S \left(   0 ,  1  \right)}{} gibt, die auf dem Rand die Identität ist.

}
{} {}

Für die folgende Aufgabe kann man
[[Obere Halbebene/Einheitskreis/Rational isomorph/Ausdehnung/Aufgabe|Aufgabe 21.8]]
heranziehen.

\inputaufgabe
{}
{

Wir betrachten die
\definitionsverweis {Mannigfaltigkeit mit Rand}{}{}

\mavergleichskettedisp
{\vergleichskette
{ M
}
{ =} { B  \left(  0 , 1 \right) \setminus \{  \left(  1   , \,   0                   \right) \}
}
{ \subseteq} { \R^2
}
{ } {
}
{ } {
}
}
{}{}{.}
Zeige, dass es eine
\definitionsverweis {stetig differenzierbare Abbildung}{}{}
\maabbdisp {\varphi} { M } { \partial M
} {}
gibt, die auf dem Rand die Identität ist.

}
{} {}

\inputaufgabe
{}
{

Zeige, dass es auf einem
\definitionsverweis {Annulus}{}{}
\definitionsverweis {bijektive}{}{}
\definitionsverweis {stetig differenzierbare}{}{} Abbildungen ohne
\definitionsverweis {Fixpunkt}{}{}
gibt.

}
{} {}

\inputaufgabe
{}
{

Zeige, dass die im Beweis zum
[[Brouwersche Fixpunktsatz/Stetig differenzierbar/Retraktion/Fakt|Brouwerschen Fixpunktsatz]]
verwendete Abbildung

\mavergleichskettedisp
{\vergleichskette
{ \varphi(x)
}
{ =} { x + \left( - \left\langle  x  , h(x)  \right\rangle + \sqrt{1 + \left\langle  x  , h(x)  \right\rangle^2 - \Vert { x} \Vert^2  }\right) \cdot h(x)
}
{ } {
}
{ } {
}
{ } {
}
}
{}{}{}
\zusatzklammer {mit

\mavergleichskette
{\vergleichskette
{ h(x)
}
{ = }{ { \frac{   \psi(x)-x }{  \Vert { \psi (x)-x} \Vert  } }
}
{  }{
}
{    }{
}
{   }{
}
}
{}{}{}} {} {}
die Einheitskugel auf die Einheitssphäre abbildet.

}
{} {}

  \zwischenueberschrift{Aufgaben zum Abgeben}

\inputaufgabe
{4}
{

Es sei $D$ das durch
$(0,2),\, (1,-1)$ und $(-2,-1)$ gegebene Dreieck und
\mathl{\tau =( 3x^2y^5-x \sin  y ) dx \wedge dy}{} eine
$2$-\definitionsverweis {Differentialform}{}{} auf $D$. Finde eine
\definitionsverweis {Stammform}{}{} für $\tau$ und berechne damit
\mathl{\int_{ D  }  \tau}{} durch ein Integral über dem Dreiecksrand.

}
{} {}

\inputaufgabe
{6}
{

Wir betrachten den Würfel

\mavergleichskettedisp
{\vergleichskette
{ Q
}
{ =} { [-1,1]^3
}
{ \subseteq} { \R^3
}
{ } {
}
{ } {
}
}
{}{}{}
und die
$2$-\definitionsverweis {Differentialform}{}{}

\mavergleichskettedisp
{\vergleichskette
{ \omega
}
{ =} { dx \wedge dy +y dx \wedge dz +x^2y^2z^2 dy \wedge dz
}
{ } {
}
{ } {
}
{ } {
}
}
{}{}{.}
Berechne
\mathl{d \omega}{} und die beiden Integrale
\mathkor {} {\int_{  \partial Q  }   \omega} {und} {\int_{  Q  }  d \omega} {}
\zusatzklammer {unabhängig voneinander} {} {.}

}
{} {}

\inputaufgabe
{4}
{

Berechne den Flächeninhalt der
\definitionsverweis {abgeschlossenen Einheitskreisscheibe}{}{}
über ein geeignetes
\definitionsverweis {Wegintegral}{}{.}

}
{} {}

\inputaufgabe
{2}
{

Zeige, dass es auf einem
\definitionsverweis {Torus}{}{}
\definitionsverweis {bijektive}{}{}
\definitionsverweis {stetig differenzierbare}{}{} Abbildungen ohne
\definitionsverweis {Fixpunkt}{}{}
gibt.

}
{} {}

\inputaufgabe
{3}
{

Es sei $B$ die
\definitionsverweis {abgeschlossene Einheitskreisscheibe}{}{}
und $K$ der obere Halbkreisbogen. Zeige, dass es eine
\definitionsverweis {differenzierbare Abbildung}{}{}
\maabbdisp {\varphi} { B } { K
} {}
gibt, deren
\definitionsverweis {Einschränkung}{}{}
auf
\mathl{K}{} die
\definitionsverweis {Identität}{}{}
ist.

}
{} {}

\inputaufgabe
{3}
{

Es seien

\mavergleichskette
{\vergleichskette
{ v,w
}
{ \in }{ \R^n
}
{  }{
}
{    }{
}
{   }{
}
}
{}{}{}
mit
\mathkor {} {\Vert { v } \Vert \leq 1} {und} {\Vert { w } \Vert = 1} {.}
Zeige, dass es ein

\mavergleichskette
{\vergleichskette
{ a
}
{ \in }{ \R
}
{  }{
}
{    }{
}
{   }{
}
}
{}{}{}
mit

\mavergleichskette
{\vergleichskette
{ \Vert {v+aw} \Vert
}
{ = }{ 1
}
{  }{
}
{    }{
}
{   }{
}
}
{}{}{}
gibt.

}
{} {}

\inputaufgabe
{4}
{

Wir betrachten die
\definitionsverweis {Mannigfaltigkeit mit Rand}{}{}

\mavergleichskettedisp
{\vergleichskette
{ M
}
{ =} { B  \left(  0 , 1 \right) \setminus \{  \left(  1   , \,   0                   \right) \}
}
{ \subseteq} { \R^2
}
{ } {
}
{ } {
}
}
{}{}{.}
Zeige, dass es eine fixpunktfreie
\definitionsverweis {stetig differenzierbare Abbildung}{}{}
\maabbdisp {\varphi} { M } { M
} {}
gibt.

}
{} {}

 [[Kategorie:Latexseite]]
